% Matched-shape random-initialization operating point.
% Source: PAPER_B_DATA.md secs. 3(d), 8.3.
\begin{table}[t]
\centering
\small
\setlength{\tabcolsep}{4pt}
\begin{tabular}{@{}lccc@{}}
\toprule
\textbf{Arm} & \textbf{Peak LR} & \textbf{PPL}$\downarrow$ & \textbf{MMLU} \\
\midrule
inherited, train all & $2\mathrm{e}{-5}$ & $10.561$ & $.3191$ \\
inherited front frozen & $2\mathrm{e}{-5}$ & $12.797$ & $.2628$ \\
fully random init, train all & $1\mathrm{e}{-4}$ & $11.498$ & $.2461$ \\
\bottomrule
\end{tabular}
\caption{\textbf{Same-shape operating points at keep14+fresh2, all at 200k
steps.} Random init does not isolate initialization because its learning rate
also differs. Frozen-front is learning-rate matched to train-all but changes the
set of trainable modules; neither comparison alone identifies a causal factor.}
\label{tab:control}
\end{table}
